% Generated draft: requires mathematical and source review.
\documentclass[11pt]{article}
\usepackage[margin=25mm]{geometry}
\usepackage{fontspec}
\setmainfont{Latin Modern Roman}
\usepackage{amsmath,amssymb,mathtools,unicode-math}
\setmathfont{Latin Modern Math}
\usepackage{enumitem,needspace,xurl,hyperref}
\setlength{\parindent}{0pt}
\setlength{\parskip}{5pt}
\setlength{\emergencystretch}{3em}
\newcommand{\N}{\mathbb{N}}
\newcommand{\Z}{\mathbb{Z}}
\newcommand{\Q}{\mathbb{Q}}
\newcommand{\R}{\mathbb{R}}
\newcommand{\C}{\mathbb{C}}
\newcommand{\F}{\mathbb{F}}
\newcommand{\A}{\mathbb{A}}
\newcommand{\PP}{\mathbb{P}}
\newcommand{\E}{\mathbb{E}}
\newcommand{\eps}{\varepsilon}
\newcommand{\dd}{\,\mathrm{d}}
\begin{document}
\section*{UnsolvedMath ID: 701562}
\textbf{Source problem code:} catalogue entry
\subsection{Definitions and mathematical statement}
Work in the usual language of groups, with multiplication, inversion, and identity. A group quasi-identity is a universally quantified implication
\[
\forall \mathbf{x}\;\left(\bigwedge_{i=1}^{m} w_i(\mathbf{x})=1\ \Longrightarrow\ w(\mathbf{x})=1\right),
\]
where the expressions are group words and \(m\geq 0\); an empty conjunction is true. A group quasivariety is the class of all groups satisfying some collection of group quasi-identities. For a class \(\mathcal C\) of groups, let \(\operatorname{Qvar}(\mathcal C)\) denote the class of all groups satisfying every quasi-identity valid in every member of \(\mathcal C\). Write \(\operatorname{Qvar}(G)=\operatorname{Qvar}(\{G\})\). A subquasivariety of a group quasivariety \(\mathcal V\) is a group quasivariety \(\mathcal W\) with \(\mathcal W\subseteq\mathcal V\).

For a finite group \(G\), let \(\operatorname{Syl}(G)\) be the collection of all its Sylow \(p\)-subgroups, for all primes \(p\). Here a Sylow \(p\)-subgroup has order equal to the largest power of \(p\) dividing \(|G|\), allowing order \(1\). In particular, \(\operatorname{Syl}(1)=\{1\}\). The notation \(\operatorname{Qvar}(\operatorname{Syl}(G))\) denotes the quasivariety generated jointly by these subgroups.

\textbf{Question (A. I. Budkin; Kourovka Notebook 19.14).} For every finite group \(G\), does the following implication hold?
\[
\begin{gathered}
\operatorname{Qvar}(\operatorname{Syl}(G))\text{ has only finitely many subquasivarieties}
\\[2pt]
\Longrightarrow
\\[2pt]
\operatorname{Qvar}(G)\text{ has only finitely many subquasivarieties}.
\end{gathered}
\]
The quasivarieties and their subquasivarieties are understood in the conventional sense as classes of arbitrary groups, although the generating group \(G\) is finite. The supplied record gives no separate formal definition of quasivariety. Its designation of this question as open is supplied registry information, not an independently verified mathematical or literature assessment.
\subsection{Short English statement}
For a finite group \(G\), suppose the quasivariety generated jointly by all Sylow subgroups of \(G\) has finitely many subquasivarieties. Must the quasivariety generated by \(G\) also have finitely many subquasivarieties? The question is attributed to A. I. Budkin.
\subsection{Sources}
\begin{itemize}
\item Supplied local UnsolvedMath registry record, permanent ID \(701562\), exact source code \texttt{catalogue entry}, titled \emph{Kourovka Notebook Problem 19.14}. The record attributes the problem to A. I. Budkin, gives the proposed year as 2018, and labels it open. Its supplied update date is \texttt{2026-10-16}; this is registry metadata, not independent source verification.
\item \emph{The Kourovka Notebook (Unsolved Problems in Group Theory)}, original problem 19.14, attributed in the supplied record to the 19th issue (2018) and the 21st issue edition (2026), \texttt{arXiv:1401.0300v48}. The supplied collection description names E. I. Khukhro and V. D. Mazurov as editors. These bibliographic details are reproduced from the supplied record; the cited editions have not been independently consulted.
\end{itemize}
\end{document}
